Throughout, $(v, w)$ is a fixed switch in $E$, for instance set and clear in a basis of $F^{2}$ (Theorem \claimref{synthesis.set-and-clear-of-a-bit-generate-the-hyperbolic-plane}(i)), $h = wv - vw$, and $hw = w$, $wh = -w$, $hv = -v$, $vh = v$, $h^{2} = 1$ by the relations of a switch.

\emph{(i)} Every $x \in E$ satisfies $x^{2} = (\operatorname{tr}x)\,x - \det x$, so $x^{2}$ is a scalar exactly when $x$ is a scalar or $\operatorname{tr}x = 0$. Suppose $x \in V$ has $\operatorname{tr}x \ne 0$. Then $x$ is a scalar $\lambda \ne 0$. Take $y \in V$ independent of $x$; $y$ is not a scalar, so $\operatorname{tr}y = 0$, and $x + y$ is not a scalar and has trace $2\lambda \ne 0$, so $(x + y)^{2}$ is not a scalar, a contradiction. On $E_{0}$, $x^{2} = -\det x$. By Theorem \claimref{synthesis.set-and-clear-of-a-bit-generate-the-hyperbolic-plane}(ii) and (iv), $h$, $v + w$, $v - w$ is a basis of $E_{0}$ of pairwise anticommuting elements of squares $1$, $1$, $-1$, so $Q(\alpha h + \sigma(v + w) + \tau(v - w)) = \alpha^{2} + \sigma^{2} - \tau^{2}$.

\emph{(ii)} A nonzero $n \in E_{0}$ with $n^{2} = 0$ is a null vector of $Q$, and for $x, y \in E_{0}$, $xy + yx = (x + y)^{2} - x^{2} - y^{2}$ is the scalar $2B(x,y)$, $B$ the polar form of $Q$. Let $n_{1}, n_{2} \in E_{0}$ be independent with $n_{1}^{2} = n_{2}^{2} = 0$. A nonzero map of square zero on $F^{2}$ is conjugate to $w$, so choose the basis of $F^{2}$ with $n_{1} = w$. For $y = \alpha h + \beta w + \gamma v$, $wy + yw = \alpha(wh + hw) + \gamma(wv + vw) = \gamma$. If $n_{1}n_{2} + n_{2}n_{1} = 0$, then $n_{2} = \alpha h + \beta w$ and $n_{2}^{2} = \alpha^{2} + \alpha\beta(hw + wh) = \alpha^{2} = 0$, so $n_{2}$ is a multiple of $w$, which is excluded. Hence $n_{1}n_{2} + n_{2}n_{1} = \gamma \ne 0$. In particular $(n_{1} + n_{2})^{2} = \gamma \ne 0$: no plane is totally null. And $(n_{1}, \gamma^{-1}n_{2})$ is a switch spanning the plane. Conversely, on a plane spanned by a switch $(a, b)$, $Q(\alpha a + \beta b) = \alpha\beta$, with exactly the two null lines $Fa$ and $Fb$. If $Q|_{V}$ is degenerate, its radical is a line (not $V$, as $V$ is not totally null) of null vectors, and it is the only null line, since a second would be orthogonal to it and make $V$ totally null. If $Q|_{V}$ is nondegenerate and has a null vector $n$, choose $y \in V$ with $ny + yn \ne 0$; then $y - \frac{Q(y)}{ny + yn}\,n$ is null (expand the square) and independent of $n$. So a nondegenerate plane has no null line or two.

\emph{(iii)} Let $x, y$ be an orthogonal basis of $V$ (it exists since $2 \ne 0$). The plane $P = \operatorname{span}\{h, v + w\}$ has $Q > 0$ on its nonzero vectors, a sum of two squares in an ordered field. $V \cap P \ne 0$ by dimensions, so $Q$ takes a positive value on $V$, and not both of $Q(x), Q(y)$ are $\le 0$. If both were negative, $Q$ would be negative on the nonzero vectors of $V$, against $V \cap P \ne 0$. If one were negative and the other zero, or both zero, the same contradiction or (ii) applies ($Q \le 0$ on $V$, respectively $V$ totally null). This leaves $(+,+)$, $(+,-)$, $(+,0)$. A switch plane contains $a + b$ and $a - b$ of squares $1$ and $-1$, orthogonal: case $(+,-)$. Conversely if $Q(x) = \lambda > 0 > -\mu = Q(y)$ and $\mu/\lambda = \rho^{2}$, then $\rho x \pm y$ are independent null vectors, and (ii) applies.

\emph{(iv)} If $Q|_{V}$ is nondegenerate, take an orthogonal basis $x, y$: they anticommute and have nonzero scalar squares, so both are invertible. If $\alpha + \beta x + \gamma y + \delta xy = 0$, conjugating by $x$ changes the sign of the last two terms, so (as $2 \ne 0$) $\alpha + \beta x = 0$ and $(\gamma + \delta x)y = 0$, hence $\gamma + \delta x = 0$; since $x$ is not a scalar, all four coefficients vanish. So $1, x, y, xy$ are independent and span $E$. If $Q|_{V}$ is degenerate, take $n$ in the radical, a basis of $F^{2}$ with $n = w$, and $y \in V$ independent of $n$; by the computation in (ii), $y = \alpha h + \beta w$ with $\alpha \ne 0$, and $yw = \alpha w$, $wy = -\alpha w$, so the algebra generated is $\operatorname{span}\{1, w, h\}$, the upper triangular matrices in the basis $(s_{\mathrm{off}}, s_{\mathrm{on}})$.

\emph{(v)} By (ii) any switch $(a, b)$ spanning $V$ has $a$ and $b$ on the two null lines of $V$, so $\{\ker a, \ker b\}$ depends only on $V$. A nonzero map of square zero on $F^{2}$ has image equal to its kernel, a line. If $\ker a = \ker b = \ell$ then $\operatorname{im}a = \operatorname{im}b = \ell$ and $ab = ba = 0$, against $ab + ba = 1$: the two lines differ. A nonzero map of square zero is determined up to a scalar by its kernel, so the pair of lines determines the two null lines of $V$ and hence $V$: the map is injective. Given distinct lines $\ell_{0}, \ell_{1}$, choose $s_{\mathrm{off}} \in \ell_{0}$, $s_{\mathrm{on}} \in \ell_{1}$; set and clear in this basis are a switch with $\ker w = \ell_{0}$, $\ker v = \ell_{1}$: the map is onto. $\mathrm{GL}_{2}(F)$ acts transitively on pairs of distinct lines, and conjugation carries switch planes to switch planes compatibly with the map.
